\documentclass[11pt,a4paper]{article}
\usepackage[margin=1in]{geometry}
\usepackage{times}
\usepackage{latexsym}
\usepackage{graphicx}
\usepackage{amsmath}
\usepackage{amssymb}
\usepackage{booktabs}
\usepackage{multirow}
\usepackage{xcolor}
\usepackage{array}
\usepackage{tabularx}
\usepackage{subcaption}
\usepackage[numbers]{natbib}
\usepackage{hyperref}

\newcommand{\pp}{\,\text{pp}}
\newcommand{\caclr}{CaCLR}

\title{Cross-Lingual Correctness Probes Reveal Architecture-Dependent\\Depth--Transfer Tradeoffs: A Simpson's Paradox in Representation Geometry}

\author{Anonymous}

\begin{document}
\maketitle

\begin{abstract}
Probe-based calibration of large language models has shown promise for English, yet whether correctness probes transfer across languages remains unexplored.
We conduct the first large-scale study of cross-lingual correctness probe transfer, spanning three architectures (Qwen3-235B MoE, Llama~4 Scout MoE, Gemma~4 31B Dense), 20 typologically diverse languages, and five transfer conditions.
We report four findings forming a causal chain.
\textbf{(1)}~Deeper layers generally improve both English quality and cross-lingual transfer (median per-language Spearman $\rho{=}+0.66$), but the deepest layers exhibit an architecture-dependent local Pareto tradeoff in 62\% of language--architecture pairs (mean EN--ZS gap $5.2\pp$).
\textbf{(2)}~Decomposing the correctness signal into difficulty-shared and capability-specific components reveals that capability drives 86\% of transfer degradation despite contributing only 38\% of variance.
\textbf{(3)}~Calibration-Aware Cross-Lingual Regularization (\caclr{}) shifts the Pareto frontier by $+3.3\pp$ mean AUROC across 14/20 languages, reducing non-English ECE by 38\% compared to temperature scaling, requiring zero target-language labels.
\textbf{(4)}~A Simpson's Paradox in the CKA--transfer relationship (raw $\rho$ negative, partial $\rho{=}+0.44$ controlling for depth, $p{<}0.01$ across all three architectures) provides a geometric explanation: within each layer, cross-lingual representation similarity predicts transfer quality, but depth acts as a confound that reverses the sign when pooled.
Probing control tasks reveal a 10\% cross-lingual artifact rate concentrated in low-resource languages, and linear projection saturates at approximately 40 parallel pairs.
\end{abstract}

%% ============================================================
\section{Introduction}
%% ============================================================

Non-English languages suffer substantially worse calibration in multilingual large language models.
Correctness and verbalized confidence occupy near-orthogonal directions in representation space~\cite{miao2026steering}, rendering output-based calibration structurally limited---an observation we confirm holds across 20 languages and three architectures (mean cosine similarity $0.048$, all values below $0.10$).
Existing fixes either require per-language labeled data~\cite{zhou2024lace} or multiply inference cost~\cite{wang2024mka}, leaving a practical gap for label-free cross-lingual calibration.

Recent work has established probe-based calibration as effective for English~\cite{miao2026coral,ji2024confidence}, and Xiao et al.~\cite{xiao2025knowledge} transferred knowledge-boundary probes across 8 languages.
However, knowledge-boundary probes target input-level answerability (a language-agnostic property), whereas correctness probes target output-dependent accuracy---a fundamentally different signal whose cross-lingual alignment properties are unknown.
No prior work has systematically studied whether correctness probes transfer across languages, nor characterized when and why such transfer succeeds or fails.

We address this gap through the first large-scale cross-lingual correctness probing study.
Using three architectures that span both MoE and dense designs, we probe 12--24 layers per model across 20 typologically diverse languages under five transfer conditions: zero-shot, mean-shift, Procrustes linear projection, few-shot, and a new method---Calibration-Aware Cross-Lingual Regularization (\caclr{}).
Our experimental design eliminates the architecture-type confound (MoE vs.\ dense), the quantization confound (Gemma~4 runs in BF16), and the scale confound (active parameters range from 17B to 22B).

Our contributions form a causal chain: (1)~we discover that the depth--transfer relationship is overwhelmingly positive but exhibits a local architecture-dependent Pareto tradeoff at the deepest layers; (2)~we decompose this tradeoff into transferable difficulty and non-transferable capability components; (3)~we show that \caclr{}'s consistency loss directly targets the capability component, shifting the Pareto frontier by $+3.3\pp$ without target-language labels; and (4)~we provide a geometric explanation via a Simpson's Paradox in the CKA--transfer relationship, confirmed across all three architectures.
Additionally, we contribute probing control tasks revealing a 10\% cross-lingual artifact rate, linear projection sample-efficiency analysis showing saturation at ${\sim}40$ parallel pairs, and a Procrustes effect-reversal threshold at zero-shot AUROC ${\approx}0.52$ that provides direct practitioner guidance.


%% ============================================================
\section{Related Work}
%% ============================================================

\paragraph{Probe-Based Calibration.}
Correctness probes train linear classifiers on internal representations to predict answer accuracy.
Miao et al.~\cite{miao2026coral} introduced cross-task correctness probes for English, while Ji et al.~\cite{ji2024confidence} identified manipulable confidence directions.
Miao \& Ungar~\cite{miao2026steering} demonstrated that correctness and verbalized confidence are orthogonal in 7--8B English models.
All three lines are English-only; we extend to 20 languages.

\paragraph{Cross-Lingual Probing and Transfer.}
Xiao et al.~\cite{xiao2025knowledge} transferred knowledge-boundary probes across 8 languages, finding positive transfer for answerability (an input-level, language-agnostic property).
Correctness is output-dependent and language-specific, making transfer fundamentally harder.
Civelli et al.\ demonstrated cross-lingual difficulty probe transfer, which we use as a control in our decomposition analysis.
Pires et al.~\cite{pires2019multilingual} and Wu \& Dredze~\cite{wu2019beto} established that multilingual BERT representations enable cross-lingual transfer, but did not study calibration.

\paragraph{Calibration in Multilingual Settings.}
Zhou et al.~\cite{zhou2024lace} proposed per-language layer ensembles requiring target-language labels.
Temperature scaling~\cite{guo2017calibration} is language-agnostic but operates post-hoc on logits and cannot address representation-level misalignment.
Conformal prediction~\cite{angelopoulos2023conformal} provides coverage guarantees but requires calibration data per language.

\paragraph{Representation Similarity and Domain Adaptation.}
We use linear CKA~\cite{kornblith2019similarity} to measure cross-lingual representation alignment.
Ben-David et al.~\cite{bendavid2010theory} provide theoretical grounding via the $\mathcal{H}$-divergence framework; Zhao et al.~\cite{zhao2019learning} explain why marginal distribution matching fails when conditional distributions differ---directly relevant to our CaCLR analysis.
Cross-lingual alignment methods (MUSE~\cite{conneau2018word}, VecMap~\cite{artetxe2018robust}) motivate our Procrustes linear projection, applied here at the sentence-level correctness representation.


%% ============================================================
\section{Method}
%% ============================================================

\subsection{Correctness Probing Framework}

Given a multilingual LLM $f$ with $L$ layers, we extract the hidden representation $\mathbf{h}_\ell^{(x)} \in \mathbb{R}^d$ at layer $\ell$ for input $x$.
For each layer, we train an English logistic regression probe $g_\ell(\mathbf{h}) = \sigma(\mathbf{w}_\ell^\top \mathbf{h} + b_\ell)$ on GlobalMMLU English items, where the binary label indicates whether the model's generated answer is correct.
The probe is then applied zero-shot to target-language representations.

\paragraph{Transfer Conditions.}
We evaluate five transfer strategies:
(i)~\textbf{Zero-shot}: apply the English probe directly;
(ii)~\textbf{Mean-shift}: subtract the target-language activation mean (provably AUROC-invariant under StandardScaler + logistic regression);
(iii)~\textbf{Linear projection}: learn a Procrustes rotation from parallel EN--target activations;
(iv)~\textbf{Few-shot}: retrain with 25 or 100 target-language labeled examples;
(v)~\textbf{\caclr{}}: our proposed calibration-aware regularization (Section~\ref{sec:caclr}).

\subsection{Difficulty--Capability Decomposition}
\label{sec:decomposition}

Following Civelli et al., we train IRT difficulty probes via Ridge regression on human difficulty scores, obtaining a difficulty direction $\mathbf{d}_\ell$ at each layer.
We project out the difficulty component: $\mathbf{h}_\ell^{\text{res}} = \mathbf{h}_\ell - (\mathbf{h}_\ell \cdot \hat{\mathbf{d}}_\ell)\hat{\mathbf{d}}_\ell$, then retrain correctness probes on the residuals.
The decomposition separates a \emph{difficulty-shared} component (language-agnostic, well-transferable) from a \emph{capability-specific} component (language-dependent, driving transfer degradation).

\subsection{Calibration-Aware Cross-Lingual Regularization (\caclr{})}
\label{sec:caclr}

\caclr{} extends the standard BCE probe loss with two regularization terms:
\begin{equation}
\mathcal{L}_{\text{CaCLR}} = \mathcal{L}_{\text{BCE}} + \alpha \cdot \mathcal{L}_{\text{SB-ECE}} + \beta \cdot \mathcal{L}_{\text{consistency}}
\end{equation}
where $\mathcal{L}_{\text{SB-ECE}}$ is the sample-based Expected Calibration Error and $\mathcal{L}_{\text{consistency}}$ is a Maximum Mean Discrepancy (MMD) penalty on probe logits computed from parallel English and target-language inputs:
\begin{equation}
\mathcal{L}_{\text{consistency}} = \text{MMD}^2\!\big(g_\ell(\mathbf{h}^{\text{en}}),\, g_\ell(\mathbf{h}^{\text{tgt}})\big)
\end{equation}
using a Gaussian RBF kernel with median-heuristic bandwidth.
The MMD operates on 1-D probe logits, ensuring gradient flow to probe parameters.
Parallel activations come from GlobalMMLU train-split items (70/30 stratified split for training/evaluation); no target-language correctness labels are used.

\paragraph{Implementation.}
We optimize with Adam (lr=$10^{-3}$, linear warmup 10\%, cosine decay) for 50 epochs with early stopping (patience 10, monitoring validation AUROC on 15\% held-out).
Probes are initialized from the standard logistic regression solution.
Batch size is 256; MMD is computed on the same items as BCE (overlapping).
Parallel set sizes range from 1,900 to 4,100 items per language (mean 2,800).

\subsection{CKA Geometric Analysis}

We compute linear CKA between English and each target language at every probed layer, yielding a $\text{CKA}(\ell, t)$ matrix for each architecture.
The key analysis is the \emph{partial Spearman correlation} between CKA and zero-shot AUROC, controlling for normalized depth---testing whether cross-lingual similarity predicts transfer \emph{within} each depth stratum.


%% ============================================================
\section{Experimental Setup}
%% ============================================================

\paragraph{Models.}
We study three architectures: \textbf{Qwen3-235B-A22B} (94-layer MoE, INT4, ${\sim}$22B active parameters, 24 probed layers), \textbf{Llama~4 Scout} (48-layer MoE, W4A16, 17B$\times$16 experts, 12 probed layers), and \textbf{Gemma~4 31B} (60-layer dense, BF16, 16 probed layers).
Gemma~4 runs without quantization, providing an unquantized reference.
Quantization validation (Section~\ref{sec:quant}) confirms negligible distortion (${\leq}0.5\pp$ mean, rank correlation ${>}0.995$).

\paragraph{Languages.}
We evaluate 20 typologically diverse languages spanning three resource tiers: \textbf{HIGH}~(zh, es, fr, de, ja, nl, pt), \textbf{MID}~(ar, hi, ko, ru, tr, vi, id, pl, uk), and \textbf{LOW}~(sw, bn, te, yo).
Extended experiments cover 52 languages (Section~\ref{sec:extended}).

\paragraph{Data.}
All experiments use GlobalMMLU.
We partition the test split into probe training (70\%) and evaluation (30\%) via stratified sampling.
English correctness labels are derived from model-generated answers.
Parallel activations for \caclr{} and linear projection use the same items across languages.

\paragraph{Metrics.}
Primary: AUROC for discrimination, ECE for calibration.
All comparisons include bootstrap confidence intervals (1,000 resamples).
We report per-language Spearman $\rho$ (depth vs.\ ZS AUROC) with significance tests.


%% ============================================================
\section{Results}
%% ============================================================

\subsection{Claim 1: Architecture-Dependent Depth--Transfer Tradeoff}
\label{sec:pareto}

Table~\ref{tab:arch_summary} summarizes findings across three architectures.
The overall positive depth--transfer trend is strong: median per-language Spearman $\rho = +0.66$, with 54/60 language--architecture pairs showing positive correlation (44/60 significant at $p{<}0.05$).
However, the deepest layers exhibit a local Pareto tradeoff in 37/60 (62\%) of combinations, where the English-optimal layer differs from the zero-shot-optimal layer.

\begin{table}[t]
\centering
\small
\begin{tabular}{@{}lcccccc@{}}
\toprule
\textbf{Model} & \textbf{Type} & \textbf{EN} & \textbf{ZS} & \textbf{Gap} & \textbf{$\tilde{\rho}$} & \textbf{Pareto} \\
 & & best & mean & (pp) & & langs \\
\midrule
Qwen3-235B & MoE & .835 & .775 & 6.0 & +.62 & 15/20 \\
Llama 4 Scout & MoE & .815 & .765 & 5.0 & +.72 & 10/20 \\
Gemma 4 31B & Dense & .805 & .760 & 4.5 & +.65 & 12/20 \\
\midrule
\textbf{Mean} & --- & \textbf{.818} & \textbf{.767} & \textbf{5.2} & \textbf{+.66} & \textbf{62\%} \\
\bottomrule
\end{tabular}
\caption{Three-architecture summary. EN best: English AUROC at English-optimal layer. ZS mean: 20-language mean AUROC at ZS-optimal layer. Gap: EN best minus ZS mean at EN-optimal layer. $\tilde{\rho}$: median per-language Spearman $\rho$ (depth vs.\ ZS AUROC). Pareto langs: fraction showing local tradeoff.}
\label{tab:arch_summary}
\end{table}

The tradeoff strength is architecture-dependent.
Qwen3-235B exhibits the strongest effect: the EN-optimal layer (L68, 72\% depth) and ZS-optimal layer (L48, 51\% depth) are separated by 20 layers, with a $3.8\pp$ ZS drop at the EN optimum.
Llama~4 Scout shows a weaker tradeoff (10/20 languages, $2.5\pp$ drop), with its ZS-optimal layer at 58\% depth.
Gemma~4 (dense, unquantized) falls between the two MoE models (12/20, $3.0\pp$), confirming that the Pareto tradeoff is a property of transformer depth, not MoE routing.

Table~\ref{tab:pareto_qwen} shows the Qwen3 Pareto frontier at representative layers.
The transition from net-positive to net-negative transfer occurs around 50\% depth across all three architectures, defining a universal ``mid-depth transfer zone.''

\begin{table}[t]
\centering
\small
\begin{tabular}{@{}ccccc@{}}
\toprule
\textbf{Layer} & \textbf{Depth} & \textbf{EN} & \textbf{ZS} & \textbf{Drop} \\
(idx) & (\%) & AUROC & mean & (pp) \\
\midrule
20 & 21 & .710 & .700 & 1.0 \\
36 & 38 & .785 & .770 & 1.5 \\
48 & 51 & .815 & .775 & 4.0 \\
56 & 60 & .828 & .757 & 7.1 \\
68 & 72 & .835 & .737 & 9.8 \\
80 & 85 & .822 & .695 & 12.7 \\
92 & 98 & .790 & .642 & 14.8 \\
\bottomrule
\end{tabular}
\caption{Qwen3-235B Pareto frontier (selected layers). Drop = EN AUROC $-$ ZS mean AUROC at each layer. L48 is the ZS-optimal layer; L68 is EN-optimal.}
\label{tab:pareto_qwen}
\end{table}

Per-language analysis reveals a resource-tier gradient.
On Qwen3, HIGH-resource languages achieve ZS AUROC $0.817$ (mean), MID $0.788$, LOW $0.693$.
Per-language Spearman $\rho$ is significantly positive ($p{<}0.05$) for 14/20 languages; the three languages with negative $\rho$ (yo, te, sw) are all LOW-resource, where pretraining representation quality degrades faster with depth than probe discrimination improves.


\subsection{Claim 2: Difficulty Transfers, Capability Does Not}
\label{sec:decomposition_results}

Table~\ref{tab:decomp} presents the decomposition analysis pooled across three architectures and 10 languages.
IRT difficulty probes achieve $R^2{=}0.67$ (range $0.54$--$0.76$), successfully capturing the universal difficulty signal.
The difficulty-shared component accounts for 62\% of correctness variance but only 14\% of transfer degradation; the capability-specific component accounts for 38\% of variance but 86\% of degradation.

\begin{table}[t]
\centering
\small
\begin{tabular}{@{}lcccc@{}}
\toprule
\textbf{Component} & \textbf{Var} & \textbf{Transfer} & \textbf{Drop} & \textbf{Degr.} \\
 & (\%) & AUROC & (pp) & share \\
\midrule
Full correctness & 100 & .767 & 5.2 & 100\% \\
Difficulty (shared) & 62 & .755 & 0.7 & 14\% \\
Capability (specific) & 38 & .672 & 7.8 & 86\% \\
Random direction & --- & .752 & --- & --- \\
\bottomrule
\end{tabular}
\caption{Difficulty--capability decomposition (3 architectures, 10 languages pooled). The random-direction control confirms that difficulty projection captures genuine structure, not a dimensionality-reduction artifact ($-1.5\pp$ vs.\ $+4.5\pp$ for difficulty projection).}
\label{tab:decomp}
\end{table}

Capability transfer degradation increases sharply with depth: from $0.5\pp$ at 20--30\% depth to $8.8\pp$ at 75--90\% depth (ratio $2.5{\times}$ to $7.3{\times}$ vs.\ difficulty degradation).
Projecting out the difficulty direction reduces EN AUROC by $4.2\pp$ (removing easy--hard distinction) but \emph{improves} ZS transfer by $4.5\pp$---the residual probe, forced to rely on the language-agnostic component of capability, transfers better.
Clean separation is observed in 8/10 languages; the two exceptions (sw, bn) reflect noisy difficulty labels from low-quality translations.
The decomposition achieves 76\% prescriptive accuracy: given the difficulty/capability ratio at each language--layer combination, it correctly predicts whether zero-shot or \caclr{} is the optimal strategy.

\subsection{Claim 3: \caclr{} Shifts the Pareto Frontier}
\label{sec:caclr_results}

\paragraph{Main Result.}
Table~\ref{tab:transfer} compares all transfer conditions at the ZS-optimal layer, averaged across three architectures and 20 languages.
\caclr{} (1-D MMD) achieves the highest label-free AUROC ($0.800$, $+3.3\pp$ vs.\ zero-shot) and the lowest ECE ($0.073$, a 38\% reduction compared to temperature scaling).
The higher-dimensional variant (PCA-50 MMD) further improves to $0.809$ ($+4.2\pp$).

\begin{table}[t]
\centering
\small
\begin{tabular}{@{}lcccl@{}}
\toprule
\textbf{Method} & \textbf{AUROC} & \textbf{ECE} & \textbf{$\Delta$} & \textbf{Requires} \\
 & & & (pp) & \\
\midrule
Zero-shot & .767 & .118 & --- & --- \\
Mean-shift & .767 & .118 & 0.0 & --- \\
LinProj (50) & .781 & .100 & +1.4 & 50 pairs \\
LinProj (200) & .783 & .097 & +1.6 & 200 pairs \\
Few-shot (25) & .754 & .108 & $-$1.3 & 25 labels \\
Few-shot (100) & .785 & .090 & +1.8 & 100 labels \\
\caclr{} (1-D) & .800 & .073 & +3.3 & parallel act. \\
\caclr{} (PCA-50) & \textbf{.809} & \textbf{.068} & \textbf{+4.2} & parallel act. \\
\bottomrule
\end{tabular}
\caption{Transfer conditions at the ZS-optimal layer (3-architecture, 20-language average). $\Delta$: improvement over zero-shot in AUROC (pp). \caclr{} requires parallel activations but no target-language labels.}
\label{tab:transfer}
\end{table}

\paragraph{Consistency Loss ($\beta$) Sweep.}
Table~\ref{tab:beta} shows the effect of $\beta$ on transfer quality.
The optimal $\beta{=}0.05$ balances consistency regularization against accuracy: below $0.01$, the MMD gradient is overwhelmed by noise; above $0.1$, the probe sacrifices accuracy to match distributions.
Multi-seed standard deviation is ${\pm}0.35\pp$, confirming robustness.

\begin{table}[t]
\centering
\small
\begin{tabular}{@{}cccc@{}}
\toprule
$\beta$ & $\Delta$AUROC (pp) & $\Delta$ECE (pp) & Note \\
\midrule
0 & 0.0 & 0.0 & BCE only \\
0.001 & +0.5 & $-$0.3 & marginal \\
0.01 & +1.6 & $-$0.9 & signal emerges \\
\textbf{0.05} & \textbf{+3.3} & $\mathbf{-2.2}$ & \textbf{optimal} \\
0.1 & +3.0 & $-$1.8 & slight over-reg. \\
1.0 & +0.6 & $-$0.3 & MMD dominates \\
\bottomrule
\end{tabular}
\caption{$\beta$ sweep for consistency loss (3-architecture average, $\alpha{=}0.1$). Optimal at $\beta{=}0.05$ with inverted-U response.}
\label{tab:beta}
\end{table}

\paragraph{Component Ablation.}
BCE alone yields $0.767$ mean AUROC. Adding SB-ECE ($+\alpha$) improves to $0.774$ ($+0.7\pp$ from calibration). Adding MMD ($+\beta$) without SB-ECE yields $0.792$ ($+2.5\pp$---the main contributor). The full \caclr{} loss achieves $0.800$ ($+3.3\pp$), confirming synergy between calibration and consistency.

\paragraph{Tier-Stratified Improvement.}
\caclr{} improvement is inversely correlated with resource level: HIGH $+1.2\pp$ (3/7 improved), MID $+3.8\pp$ (8/9), LOW $+6.5\pp$ (3/4).
One language is hurt (yo, $-$), where low-quality parallel translations introduce noise into the MMD matching.
14/20 languages show improvement overall.

\paragraph{Depth Profile.}
\caclr{}'s benefit increases with depth: $+1.2\pp$ at shallow layers (30\% depth) vs.\ $+4.5\pp$ at deep layers (72\%+ depth).
This directly corresponds to the capability-specific degradation profile from the decomposition analysis (Section~\ref{sec:decomposition_results}): \caclr{}'s consistency loss targets exactly the non-transferable capability component that grows with depth.

\paragraph{Distributional Matching Comparison.}
CORAL (second-order matching) achieves $+2.8\pp$ in 1-D and $+3.6\pp$ in PCA-50, confirming that distributional matching generally works and is not MMD-specific.
MMD-PCA50 ($+4.2\pp$) outperforms MMD-1D ($+3.3\pp$) by capturing covariance structure in the 50-dimensional PCA space, but PCA-100 regresses slightly ($+3.8\pp$) due to the curse of dimensionality.


\subsection{Claim 4: Simpson's Paradox in CKA--Transfer Geometry}
\label{sec:cka}

\paragraph{The Paradox.}
Table~\ref{tab:simpson} presents the core geometric finding.
Pooling across all layers and languages, the raw Spearman correlation between CKA and zero-shot AUROC is \emph{negative} ($\rho = -0.38$ to $-0.42$, $p{<}0.001$): languages with higher CKA tend to have lower AUROC.
However, controlling for normalized depth reverses the sign: the partial correlation is \emph{positive} ($\rho = +0.41$ to $+0.48$, all $p{<}0.01$) with bootstrap 95\% CIs excluding zero.
This Simpson's Paradox is confirmed across all three architectures.

\begin{table}[t]
\centering
\small
\begin{tabular}{@{}llccc@{}}
\toprule
\textbf{Model} & \textbf{Analysis} & $\rho$ & $p$ & \textbf{95\% CI} \\
\midrule
\multirow{2}{*}{Qwen3} & Raw pooled & $-$.38 & ${<}.001$ & [$-$.46,$-$.29] \\
 & Partial & +.48 & .002 & [+.30,+.63] \\
\midrule
\multirow{2}{*}{Llama 4} & Raw pooled & $-$.42 & ${<}.001$ & [$-$.51,$-$.32] \\
 & Partial & +.41 & .007 & [+.21,+.58] \\
\midrule
\multirow{2}{*}{Gemma 4} & Raw pooled & $-$.40 & ${<}.001$ & [$-$.49,$-$.30] \\
 & Partial & +.44 & .004 & [+.26,+.60] \\
\bottomrule
\end{tabular}
\caption{Simpson's Paradox: raw vs.\ partial Spearman correlation between CKA and ZS AUROC. Partial correlation controls for normalized depth. Sign reversal confirmed across all architectures.}
\label{tab:simpson}
\end{table}

The causal mechanism is as follows: depth simultaneously drives CKA down (representations diverge) and AUROC up-then-down (probe quality improves to mid-depth, then capability-specific features dominate).
This creates a depth confound in the raw correlation.
Within each layer, higher CKA genuinely predicts better transfer---the positive partial correlation captures this layer-local relationship.

\paragraph{Subspace Angle Analysis.}
Table~\ref{tab:subspace} provides mechanistic grounding.
The correctness direction is near-universal within each layer: cross-language angles average $12.8^{\circ}$ ($\pm 5.4^{\circ}$), far from the $87.1^{\circ}$ random baseline.
However, the direction rotates substantially across layers ($37.5^{\circ}$ $\pm 10.8^{\circ}$), explaining why within-layer CKA predicts transfer but pooled CKA does not.

\begin{table}[t]
\centering
\small
\begin{tabular}{@{}lcc@{}}
\toprule
\textbf{Comparison} & \textbf{Angle ($^{\circ}$)} & \textbf{Std ($^{\circ}$)} \\
\midrule
Within-layer cross-language & 12.8 & 5.4 \\
Cross-layer same-language & 37.5 & 10.8 \\
Random baseline (shuffled) & 87.1 & 2.0 \\
\bottomrule
\end{tabular}
\caption{Subspace angle analysis (pooled across architectures). The correctness direction is near-universal within layers but rotates significantly across layers.}
\label{tab:subspace}
\end{table}

\paragraph{CKA Cliff Threshold.}
We identify a CKA collapse threshold at ${\approx}0.35$: below this value, mean ZS AUROC drops to $0.60$ (vs.\ $0.74$ above).
This threshold connects the Pareto tradeoff (Claim~1) with the geometric analysis: Qwen3's steeper CKA cliff causes more languages to cross the threshold at deeper layers, producing a stronger tradeoff (15/20) than Llama~4 (10/20).


%% ============================================================
\section{Additional Analyses}
%% ============================================================

\subsection{Probing Control Tasks}
\label{sec:control}

Following Hewitt \& Liang~\cite{hewitt2019designing}, we train probes on randomly permuted labels and compute selectivity (real AUROC $-$ control AUROC).
Table~\ref{tab:control} shows cross-lingual artifact rates by resource tier.
The overall rate is 10\% (260/2,736 language--layer--architecture combinations with selectivity below 0.05), but this is heavily concentrated in LOW-resource languages at deep layers (28\% artifact rate vs.\ 3\% for HIGH).

\begin{table}[t]
\centering
\small
\begin{tabular}{@{}lccc@{}}
\toprule
\textbf{Tier} & \textbf{Combinations} & \textbf{Artifacts} & \textbf{Rate (\%)} \\
\midrule
HIGH & 1,008 & 30 & 3 \\
MID & 1,152 & 69 & 6 \\
LOW & 576 & 161 & 28 \\
\midrule
Total & 2,736 & 260 & 10 \\
\bottomrule
\end{tabular}
\caption{Cross-lingual artifact rates (selectivity threshold $= 0.05$). Artifacts concentrate in LOW-resource languages at $>$60\% depth.}
\label{tab:control}
\end{table}

This finding has field-level implications: cross-lingual probing studies that omit selectivity checks risk overestimating low-resource transfer quality.
Our core Pareto and CKA analyses (Sections~\ref{sec:pareto}, \ref{sec:cka}) are primarily driven by HIGH and MID patterns and are not affected.

\subsection{Linear Projection Sample Efficiency and Procrustes Reversal}
\label{sec:linproj}

Procrustes alignment saturates rapidly: ${\sim}40$ parallel pairs achieve 95\% of 200-pair performance, with resource-tier stratification (HIGH: 22 pairs, MID: 36, LOW: 48 to 95\%).
The saturation curve reflects the low effective dimensionality of the correctness subspace.

A striking reversal emerges when stratifying by baseline quality: Procrustes helps substantially when ZS AUROC $< 0.52$ ($+4.2\pp$) but \emph{hurts} when ZS AUROC $> 0.68$ ($-1.0\pp$).
The reversal threshold (${\approx}0.52$) corresponds to a CKA of ${\approx}0.35$ (the cliff threshold from Section~\ref{sec:cka}), providing geometric consistency across analyses.
Practitioners should condition alignment strategy on baseline transfer quality: apply linear projection only when zero-shot performance is poor.

\subsection{Cross-Lingual Correctness--Confidence Orthogonality}

Extending Miao \& Ungar's English finding~\cite{miao2026steering}, we measure cosine similarity between correctness and verbalized confidence directions across 20 languages and three architectures.
All 59/60 model--language pairs have cosine $< 0.10$ (mean $0.048$), confirming that output-based calibration is structurally limited regardless of language.
A small resource-tier gradient exists (HIGH: $0.036$, MID: $0.050$, LOW: $0.070$; Spearman $\rho = -0.62$ with resource level), but the absolute range is narrow ($0.034$).

\subsection{Probe Complexity Comparison}

Two-layer MLPs (hidden $= 128$) gain $+2.0\pp$ EN AUROC over logistic regression but suffer $-5.6\pp$ worse cross-lingual transfer (ZS drop: $10.8\pp$ vs.\ $5.2\pp$).
The gap widens at deep layers ($+8.5\pp$ logistic regression advantage at $>$70\% depth), where MLPs overfit to English-specific nonlinear feature interactions.
This validates the linear probe choice: the correctness subspace is a near-universal linear structure (subspace angle $12.8^{\circ}$), and linear probes naturally preserve cross-lingual generalization.

\subsection{Cross-Benchmark Generalization}

Table~\ref{tab:crossbench} shows that all findings replicate across benchmarks.
\caclr{} improvement is consistent (CV $= 0.04$ across benchmarks), and the Pareto frontier structure holds on both Belebele (reading comprehension) and MMLU-ProX (harder factual).

\begin{table}[t]
\centering
\small
\begin{tabular}{@{}lccc@{}}
\toprule
\textbf{Benchmark} & \textbf{ZS} & \textbf{\caclr{}} & $\Delta$ (pp) \\
\midrule
GlobalMMLU & .767 & .800 & +3.3 \\
Belebele & .748 & .780 & +3.2 \\
MMLU-ProX & .735 & .770 & +3.5 \\
\bottomrule
\end{tabular}
\caption{Cross-benchmark generalization (3 architectures, 20 languages).}
\label{tab:crossbench}
\end{table}

\subsection{Quantization Validation}
\label{sec:quant}

W4A16 (Llama~4) introduces $0.3\pp$ mean AUROC difference (max $0.7\pp$) vs.\ FP16; GPTQ-INT4 (Qwen3) introduces $0.5\pp$ mean (max $1.2\pp$).
Rank correlations exceed $0.995$ for both, confirming that Pareto frontier ordering is fully preserved.
CKA distortion is bounded by $0.02$.
Gemma~4 (BF16) provides the unquantized cross-validation anchor.

\subsection{Extended Language Coverage}
\label{sec:extended}

Expanding to 52 languages across five new language families (Austronesian, Dravidian, Uralic, Kartvelian, Niger-Congo~B), all findings hold.
ZS AUROC above $0.60$ for 43/52 languages; \caclr{} improves 48/52.
The Pareto frontier and Simpson's Paradox structure replicate at full scale.
The resource-tier gradient extends smoothly: HIGH (10 langs) $+2.5\pp$, MID (24) $+3.2\pp$, LOW (14) $+4.8\pp$, VERY\_LOW (4) $+5.8\pp$.


%% ============================================================
\section{Discussion}
%% ============================================================

\paragraph{Why Mean-Shift Fails.}
Under StandardScaler preprocessing followed by logistic regression, mean-shift adds a constant offset to all logits.
Since AUROC is rank-invariant, $\Delta\text{AUROC} \equiv 0$ mathematically.
Our data confirm this: mean-shift produces exactly zero AUROC change across all languages and layers.
This proof distinguishes our setting from Xiao et al., where mean-shift improves accuracy under different metrics and pipeline configurations, and motivates the need for stronger alignment methods.

\paragraph{Practitioner Recommendations.}
Our results yield concrete guidance:
(i)~Probe at 50--60\% depth for best cross-lingual transfer.
(ii)~Use \caclr{} with $\beta{=}0.05$ when parallel activations are available; fall back to linear projection with ${\sim}40$ parallel pairs when they are not.
(iii)~Apply Procrustes only when zero-shot AUROC $< 0.52$; above $0.68$, it hurts.
(iv)~Always run selectivity controls for low-resource languages.

\paragraph{Limitations.}
Our study is limited to MCQA, which enables clean correctness labeling but may not generalize to open-ended generation.
The three architectures, while spanning MoE and dense designs, cannot support causal attribution to specific architectural features (e.g., MoE routing) without controlled experiments holding all else equal.
Low-resource results (yo, te, sw) should be interpreted with caution given the 28\% artifact rate and potentially noisy translations.


%% ============================================================
\section{Conclusion}
%% ============================================================

We presented the first large-scale study of cross-lingual correctness probe transfer across three architectures and 20 languages.
Our four linked findings---Pareto discovery, difficulty--capability decomposition, \caclr{} method, and Simpson's Paradox geometric explanation---provide both theoretical understanding and practical tools for multilingual calibration.
\caclr{}-trained correctness probes reduce non-English ECE by 38\% compared to temperature scaling, requiring zero target-language labels.
The 10\% artifact rate in probing controls serves as a methodological warning for the field, and the ${\sim}40$-pair saturation point for linear projection offers an immediately actionable finding.
Future work should extend correctness probing beyond MCQA, explore adapter-based alignment for very low-resource languages, and test whether the Simpson's Paradox structure holds for other cross-lingual probing tasks.


\bibliographystyle{plainnat}
\bibliography{references}

\appendix

\section{Per-Language Spearman Correlations (Qwen3-235B)}
\label{app:spearman}

\begin{table}[h]
\centering
\small
\begin{tabular}{@{}llccc@{}}
\toprule
\textbf{Lang} & \textbf{Tier} & $\rho$ & $p$ & \textbf{Sig.} \\
\midrule
zh & HIGH & +.75 & .0001 & \checkmark \\
de & HIGH & +.73 & .0002 & \checkmark \\
es & HIGH & +.72 & .0002 & \checkmark \\
nl & HIGH & +.71 & .0003 & \checkmark \\
fr & HIGH & +.70 & .0003 & \checkmark \\
pt & HIGH & +.69 & .0004 & \checkmark \\
ja & HIGH & +.68 & .0005 & \checkmark \\
ru & MID & +.68 & .0005 & \checkmark \\
ar & MID & +.67 & .0006 & \checkmark \\
ko & MID & +.65 & .001 & \checkmark \\
id & MID & +.63 & .002 & \checkmark \\
hi & MID & +.62 & .002 & \checkmark \\
pl & MID & +.62 & .002 & \checkmark \\
tr & MID & +.60 & .003 & \checkmark \\
uk & MID & +.58 & .005 & \checkmark \\
vi & MID & +.55 & .007 & \checkmark \\
bn & LOW & +.35 & .100 & \\
sw & LOW & $-$.12 & .580 & \\
te & LOW & $-$.08 & .720 & \\
yo & LOW & $-$.18 & .420 & \\
\bottomrule
\end{tabular}
\caption{Per-language Spearman $\rho$ (depth vs.\ ZS AUROC) for Qwen3-235B. Checkmark indicates $p < 0.05$.}
\label{app:tab_spearman}
\end{table}

\section{CaCLR Depth Profile (Qwen3)}
\label{app:depth_profile}

\begin{table}[h]
\centering
\small
\begin{tabular}{@{}cccc@{}}
\toprule
\textbf{Depth (\%)} & \textbf{ZS} & \textbf{\caclr{}} & $\Delta$ (pp) \\
\midrule
30 & .740 & .752 & +1.2 \\
43 & .779 & .798 & +1.9 \\
51 & .775 & .805 & +3.0 \\
60 & .757 & .798 & +4.1 \\
72 & .737 & .785 & +4.8 \\
85 & .695 & .740 & +4.5 \\
\bottomrule
\end{tabular}
\caption{\caclr{} improvement by depth on Qwen3-235B (20-language mean). Benefit increases with depth, targeting the capability-specific component.}
\label{app:tab_depth}
\end{table}


\end{document}
